\section{検証 II：2 箱}\label{sec:two}

この節の双曲結び目の入力は予想式（\S\ref{sec:inputs}）であり，主張は全て「その予想式を入力とした場合」のものである．
トーラス結び目と代数的なケーブルの入力は定義どおりの refined 不変量である．

\subsection{ツイスト結び目}
$K_p=K_{p,1}$ を KNTZ~\cite{Kameyama19} の規約のツイスト結び目とする（$K_1=\mathbf{3_1}$，$K_2=\mathbf{5_2}$，$K_{-1}=\mathbf{4_1}$，$K_{-2}=\mathbf{6_1}$．GNSSS~\cite{Gukov15} の規約と同じ）．
入力は KNTZ の refined の式である（\S\ref{sec:inputs}）．
$X=\ba^2\bt^2$，$Y=\ba^4\bq^4\bt^8$，$\bar Y=\ba^4\bq^{-4}\bt^4$ とおく．

\begin{lemma}\label{lem:twist-expsum}
$\hf_\rho(K_p)$（$|\rho|=2$）は $p$ について根 $\{1,X,X^2,Y,\bar Y\}$ の指数和である（$p<0$ では根を逆数にし，全体に $\bt^2$ が掛かる）．
\end{lemma}
\begin{proof}
KNTZ の式では，$p$ は捻りの因子 $\mathbf{Tw}^{(p)}_\lambda$（$\lambda\subseteq\yng(2)$ または $\lambda\subseteq\yng(1,1)$）に固有値の $p$ 乗としてだけ入る．
太字の変数では，これは $\bigl((\ba\bt)^{2j}(\bt\bq)^{2j(j-1)}\bigr)^p$（$j=0,1,2$）の形である（GNSSS の式でも同じ形）．
入力の根は $P_{\yng(1)}$ で $\{1,X\}$，$P_{\yng(2)}$ で $\{1,X,Y\}$，$P_{\yng(1,1)}$ で $\{1,X,\bar Y\}$ である．
$\hf$ は入力の積と Adams 作用の一次結合を，命題~\ref{prop:two-necessary} の分母で割ったものである．
\end{proof}

% [R2]
\begin{theorem}\label{thm:twist}
KNTZ の refined の式を入力とすると，全ての $p\neq0$ で $\hf_{\yng(2)}(K_p)$，$\hf_{\yng(1,1)}(K_p)$ は正である．
\end{theorem}
\begin{proof}[証明の概略]
並び $(1,X,X^2,Y,\bar Y)$ のニュートン形 $s_{1+n}=\sum_{k=0}^{4}c_k\,h_{n-k}(m_0,\dots,m_k)$ を作る．
$c_k$ は $s_1,\dots,s_5$（$|p|=1,\dots,5$ の値）の整数係数の一次結合なので，Laurent 多項式である．
両方の $\rho$，両方の $p$ の符号で，$c_0,\dots,c_4$ の全係数が同じ符号であることを有限個の厳密な計算で確かめる．
$h$ の係数は非負なので，全ての $p$ で展開の前が一様な符号になる．
$\sigma$ は根の集合を保つ（$Y\leftrightarrow\bar Y$）ので，$\sigma$ 不変性は $|p|\le5$ から全ての $p$ に従い，補題~\ref{lem:G} で展開の後も非負になる．
$c_2$ の正値性は構造からは導けていない．
しかし，根の部分集合ごとの完全斉次対称式の非負結合（線形計画法による証明書，有理数で厳密に検証）でも同じ結論が得られる．
この証明は $c_2$ を使わない．
\end{proof}

\begin{compresult}
$|p|\le30$ の 120 件で，spin 表示と genus 表示の両方で負の項は 0 である．
閉じた形は，作るのに使っていない $|p|=6,\dots,12,15,20,30$ で直接計算と全係数一致する．
\end{compresult}

\begin{remark}[単峰性]
$v=\bq^2\bt$ とおくと，$c_3$ と $c_4$ は構造から正であり，$\hf(K_p)$ は全ての $p$ で $v$ についての $\mathfrak{sl}_2$ 指標の正係数の和（単峰）になる（$c_0,\dots,c_2$ は有限個の判定）．
$\yng(1)$ の単峰性は Gorsky--Hogancamp--Mellit~\cite{Gorsky21} の hard Lefschetz の性質から全ての結び目で従う．
2 箱の単峰性の構造的な理由は未解明である．
\end{remark}

\subsection{ダブルツイスト結び目}
% [R4]
\begin{theorem}\label{thm:pretzel}
KNTZ の refined の式を入力とすると，ダブルツイスト結び目 $K_{m,n}$ の 2 箱は次の範囲で正である．
\begin{enumerate}
\item $m,n$ が同符号の全ての $(m,n)$．
\item $m<0<n$ で $\min(|m|,n)\le20$ の全ての $(m,n)$．
\end{enumerate}
\end{theorem}
\begin{proof}[証明の概略]
$\hf(K_{m,n})=\bt^{2(1-\mathrm{sgn})}F(m,n)$ と書くと，$F$ は $\mathbb{Z}^2$ 全体で一つの指数和である．
同符号の領域では，2 変数のニュートン形の係数が全て同じ符号になる．
異符号の領域では 2 変数のニュートン形は一様にならない．そこで直線ごとの 1 変数のニュートン形で $\min(|m|,n)\le20$ を示す．
\end{proof}

\begin{compresult}
$|m|,|n|\le20$ の 820 点と，$21\le|m|\le30$・$21\le n\le40$ の 200 点で，両方の $\rho$ についてspin 表示と genus 表示の両方で負の項は 0 である．
\end{compresult}

異符号で $|m|,n>20$ の領域は，くさび形への分割と鏡像型の恒等式で調べる．
くさびの方向のニュートン係数 $d_2,d_3,d_4$ の非正値性は，線形計画法の証明書で示した．
残りは，1 変数の三つの数列の非正値性に帰着する．三つとは，対角線 $F(-u,u)$ と，両方の $\rho$ の $d_1(u)$ である．
これらは $u\le25$ で成り立つが，一般には未証明である．

\subsection{thick な非トーラス結び目}
% [R5]
\begin{compresult}\label{cr:thick}
Stošić の quiver と KRSS の予想式から作った $\mathbf{10_{139}}$，$\mathbf{10_{152}}$ の $\yng(2)$，$\yng(1,1)$ の入力で，2 箱は正である．
次元は $441=21^2$，$841=29^2$ で，命題~\ref{prop:two-necessary} の必要条件を満たす．
$t=q$ の極限は，$(2,s)$-ケーブルから逆に解いた $\yng(2)$，$\yng(1,1)$ の彩色 HOMFLY-PT 多項式と一致する．
\end{compresult}
これらは，命題~\ref{prop:two-necessary} の必要条件を満たす thick な非トーラス結び目で 2 箱の正値性が成り立つ例である．
knots--quivers 対応の DT 不変量の正値性からは従わない．
対称表現だけの quiver 級数からは，$\yng(1,1)$ の入力を要する $\hf_{\yng(2)}$，$\hf_{\yng(1,1)}$ は出てこない．

\begin{remark}[正値性は自動的ではない]
$\mathbf{5_2}$ で $t_c$ 次数の割り当てを別の 12 通りに変えると，いずれも負の項が出る．
正値性は入力の細かい次数に依存する，非自明な性質である．
\end{remark}

\subsection{全ての framing}
% [R7]
\begin{theorem}\label{thm:two-all-x}
自明な結び目，$T_{2,3}$，$T_{2,5}$，$T_{2,7}$，$T_{3,4}$，$\mathbf{4_1}$，$K1$，$K3$ の 2 箱は，全ての $x\in\mathbb{Z}$ で正である．
$\mathbf{4_1}$ の入力は KN のデータである．
\end{theorem}
\begin{proof}
\S\ref{sec:newton} の方法による．各結び目・各 $\rho$・$x$ の偶奇と向きの四つの組で，ニュートン形の係数が一様な符号になる起点 $x_0$ がある．
起点と 0 の間の $x$ は直接判定した．
\end{proof}
負の向きの起点は，結び目が複雑になるほど 0 から遠い（$T_{2,2k+1}$ で $-(2k+2)$，$K1$ で $-20$，$K3$ で $-22$）．この意味は未解明である．

% [R8]
\begin{theorem}\label{thm:twist-x}
KNTZ の refined の式を入力とすると，ツイスト結び目の 2 箱は，$p$ と $x$ が同符号（$p>0,x\ge0$ または $p<0,x\le0$）の全ての $(p,x)$ で正である．
\end{theorem}
\begin{proof}[証明の概略]
$\bt^{-2x}\hf^{(x)}_\rho(K_p)$ は $p$ と $x$ の両方について指数和である（補題~\ref{lem:twist-expsum}，命題~\ref{prop:expsum}）．
テンソル積のニュートン形の係数を $|p|=1,\dots,5$ と 3 点の $x$ から作り，同符号の向きで一様な符号であることを確かめる．
\end{proof}

\begin{compresult}
逆符号の領域では，$|p|\le6$，$|x|\le14$ の 696 件で負の項は 0 である．この領域の証明書は得られていない．
\end{compresult}

\begin{compresult}[framing を振った直接の判定]\label{cr:framing-batch}
次の 1821 件（$(x,\rho)$ の組）で，spin 表示と genus 表示の両方で負の項は 0 である（9 件は $\hf=0$）．
3 箱までは，自明な結び目と $T_{2,3}$（$|x|\le6$），$\mathbf{4_1}$ と $T_{2,5}$（$|x|\le4$），$T_{2,7}$ と $T_{3,4}$（$|x|\le3$），
$T_{3,5},T_{2,11},T_{3,7},T_{4,5}$ と $K6a,K7,K8,K9$（$|x|\le10$），$K3$（$|x|\le2$），$K1$（$|x|\le2$ と $x=\pm6$）である．
$K10$ は 2 箱まで $|x|\le10$，3 箱は $|x|\le2$ である．
ツイスト結び目 $K_{\pm1},\dots,K_{\pm6}$（$|x|\le4$）は 2 箱までである．
最大の例は $K10$ の 3 箱，$x=2$ で，spin 表示で約 189 万項である．
\end{compresult}

\begin{table}[h]
\centering
\caption{framing を振った判定．$|\rho|=1,2,3$ の列は調べた framing $x$ の範囲，件数は判定した $(x,\rho)$ の組の数（「各」は 1 結び目あたり．$K10$ の 3 箱の $x=0$ は表~\ref{tab:three} の判定で，件数に含めない），負（spin/genus）は spin 表示と genus 表示の負の項の数の合計（全体の符号は良い方）．自明な結び目の 9 件は $\wh f=0$．}\label{tab:framing}
\resizebox{\textwidth}{!}{% scripts/paper_framing_table.py が logs/framing_batch.log と data/framed/ から生成。
% 判定件数の合計 1821（重なりは 1 件と数える）。
\begin{tabular}{llcccc}
\hline
結び目 & 箱の数: $x$ の範囲 & 件数 & 前の負 & 後の負 & $\wh f=0$ \\
\hline
自明な結び目 & 1: $-6..6$; 2: $-6..6$; 3: $-6..6$ & 78 & 0 & 0 & 9 \\
$T_{2,3}$ & 1: $-6..6$; 2: $-6..6$; 3: $-6..6$ & 78 & 0 & 0 & 0 \\
$T_{2,5}$ & 1: $-4..4$; 2: $-4..4$; 3: $-4..4$ & 54 & 0 & 0 & 0 \\
$T_{2,7}$ & 1: $-3..3$; 2: $-3..3$; 3: $-3..3$ & 42 & 0 & 0 & 0 \\
$T_{3,4}$ & 1: $-3..3$; 2: $-3..3$; 3: $-3..3$ & 42 & 0 & 0 & 0 \\
$T_{3,5}$ & 1: $-10..10$; 2: $-10..10$; 3: $-10..10$ & 126 & 0 & 0 & 0 \\
$T_{2,11}$ & 1: $-10..10$; 2: $-10..10$; 3: $-10..10$ & 126 & 0 & 0 & 0 \\
$T_{3,7}$ & 1: $-10..10$; 2: $-10..10$; 3: $-10..10$ & 126 & 0 & 0 & 0 \\
$T_{4,5}$ & 1: $-10..10$; 2: $-10..10$; 3: $-10..10$ & 126 & 0 & 0 & 0 \\
$\mathbf{4_1}$ & 1: $-4..4$; 2: $-4..4$; 3: $-4..4$ & 54 & 0 & 0 & 0 \\
$K1$ & 1: $-2..2$; 2: $-2..2$; 3: $-6,-2,-1,0,1,2,6$ & 36 & 0 & 0 & 0 \\
$K3$ & 1: $-2..2$; 2: $-2..2$; 3: $-2..2$ & 30 & 0 & 0 & 0 \\
$K6a$ & 1: $-10..10$; 2: $-10..10$; 3: $-10..10$ & 126 & 0 & 0 & 0 \\
$K7$ & 1: $-10..10$; 2: $-10..10$; 3: $-10..10$ & 126 & 0 & 0 & 0 \\
$K8$ & 1: $-10..10$; 2: $-10..10$; 3: $-10..10$ & 126 & 0 & 0 & 0 \\
$K9$ & 1: $-10..10$; 2: $-10..10$; 3: $-10..10$ & 126 & 0 & 0 & 0 \\
$K10$ & 1: $-10..10$; 2: $-10..10$; 3: $-2,-1,1,2$ & 75 & 0 & 0 & 0 \\
$K_{1}$ & 1: $-4..4$; 2: $-4..4$ & 27 & 0 & 0 & 0 \\
$K_{-1}$ & 1: $-4..4$; 2: $-4..4$ & 27 & 0 & 0 & 0 \\
$K_{2}$ & 1: $-4..4$; 2: $-4..4$ & 27 & 0 & 0 & 0 \\
$K_{-2}$ & 1: $-4..4$; 2: $-4..4$ & 27 & 0 & 0 & 0 \\
$K_{3}$ & 1: $-4..4$; 2: $-4..4$ & 27 & 0 & 0 & 0 \\
$K_{-3}$ & 1: $-4..4$; 2: $-4..4$ & 27 & 0 & 0 & 0 \\
$K_{4}$ & 1: $-4..4$; 2: $-4..4$ & 27 & 0 & 0 & 0 \\
$K_{-4}$ & 1: $-4..4$; 2: $-4..4$ & 27 & 0 & 0 & 0 \\
$K_{5}$ & 1: $-4..4$; 2: $-4..4$ & 27 & 0 & 0 & 0 \\
$K_{-5}$ & 1: $-4..4$; 2: $-4..4$ & 27 & 0 & 0 & 0 \\
$K_{6}$ & 1: $-4..4$; 2: $-4..4$ & 27 & 0 & 0 & 0 \\
$K_{-6}$ & 1: $-4..4$; 2: $-4..4$ & 27 & 0 & 0 & 0 \\
\hline
\end{tabular}
}
\end{table}
